\section{Introdução}

Os modelos de inteligência artificial (IA) estão cada vez mais presentes em sistemas do cotidiano, como telefones celulares, e em aplicações que exigem alta performance e precisão, como detecção de ameaças, sistemas militares, reconhecimento de objetos, agentes médicos, veículos autônomos e sistemas educacionais. A maioria desses modelos, entretanto, costuma prover explicações humanamente incompreensíveis para o processo de tomada de decisão. O nome \textit{black-box} é atribuído àqueles modelos que não conseguem justificar, de forma adequada ao contexto ou ao público-alvo, como chegam a determinado resultado, o que compromete a atribuição de responsabilidade, a avaliação ética e o entendimento do próprio funcionamento do sistema \citep{Ullah2024}. 

Talvez o Regulamento Geral de Proteção de Dados (GDPR) \citep{EU2016GDPRen} seja a melhor ferramenta jurídica para regular o uso de dados pessoais em IA, sendo robusto e abrangente tanto para o contexto organizacional quanto para o contexto dos sistemas automatizados \citep{Hamon2022}. O documento tem como objetivo proteger os direitos e liberdades das pessoas no tratamento de dados pessoais, assegurando simultaneamente a livre circulação desses dados na União Europeia \citep{EU2016GDPRen}. Seu artigo 5º define os princípios que regem o tratamento de dados, são eles: legalidade, lealdade, transparência, exatidão, integridade e responsabilidade. Esses princípios corroboram a confiança em um sistema automatizado e, portanto, devem ser empregados para garantir a segurança dos dados do usuário, assim como a segurança da própria organização que disponibiliza o sistema. Já os artigos 13(2)(f), 14(2)(g), 15(1)(h) e 22º garantem que os titulares dos dados sejam informados sobre a existência de processos automatizados de decisão, a lógica envolvida e as consequências potenciais desses processos, além de assegurarem o direito de não serem submetidos exclusivamente a decisões automatizadas que produzam efeitos significativos. O considerando 39 informa que o tratamento dos dados pessoais deve ser feito de forma lícita e equitativa, além de reforçar que as informações sobre o tratamento de dados devem ser de fácil acesso e compreensão para as pessoas. Essa ênfase na interpretabilidade demonstra que a capacidade de entender as decisões de sistemas automatizados não é apenas um princípio ético, mas um requisito legal essencial para o desenvolvimento de sistemas de IA, capazes de garantir responsabilidade, auditabilidade e confiança \citep{Hamon2022, EU2016GDPRen}.

Definir padrões éticos universais para o uso de IA é algo desafiador e talvez impossível, visto as variações culturais, sociais e econômicas entre os países. Levar em conta a privacidade de dados, a necessidade de consentimento para uso dos dados pessoais, a equidade e a necessidade de contornar preconceitos é eticamente fundamental, ao mesmo tempo que se mostra um desafio imenso \citep{Radanliev2025}. Alguns princípios éticos comuns foram levantados por \cite{Radanliev2025} para serem levados em conta nos sistemas de IA. De acordo com o autor, desenvolver um sistema de IA claro, equitativo, sem preconceitos e que garanta a proteção dos dados pessoais significa alcançar um equilíbrio entre os diferentes princípios éticos fundamentais, de forma que eles se reforcem mutuamente. Os princípios éticos fundamentais considerados no artigo são a justiça, para garantir que as decisões serão tomadas sem preconceito ou parcialidade, a privacidade, para garantir a proteção de dados e o consentimento dos usuários, e a transparência, para garantir a clareza com relação às decisões tomadas pelo sistema.

Segundo \cite{Radanliev2025}, essas dimensões éticas da IA não estão isoladas, mas profundamente interligadas. A interação entre justiça e transparência corresponde ao princípio da responsabilidade, pois é nesse ponto que decisões algorítmicas claras e auditáveis se tornam passíveis de responsabilização. A interação entre transparência e privacidade representa a confiança do usuário, visto que a clareza sobre o uso dos dados pessoais é essencial para que os indivíduos confiem no sistema. Já a interação entre justiça e privacidade remete às práticas de dados não discriminatórias, nas quais a proteção e o uso ético das informações pessoais devem impedir que vieses e desigualdades sejam reproduzidos. Quando os três princípios se intersectam, é formado o ideal de uso responsável da IA, que busca harmonizar a clareza dos modelos com a equidade nas decisões e a proteção dos dados sensíveis. Assim, as intersecções revelam que nenhum princípio ético é suficiente isoladamente: apenas sua integração conduz a um ecossistema de IA realmente confiável, justo e socialmente responsável.

O \textit{Information Commissioner’s Office} (ICO), em parceria com o \textit{Alan Turing Institute}, propõe uma estrutura organizacional que trata a explicabilidade como dever institucional e não apenas como característica técnica do modelo. Essa orientação foi desenvolvida para melhorar transparência e responsabilização no uso de IA em decisões que afetam indivíduos, especialmente no tratamento de dados pessoais, e tem como foco estabelecer boas práticas de prestação de contas, e não apenas requisitos de desempenho algorítmico. A proposta parte do princípio de que explicar uma decisão automatizada exige processos internos claros (governança, papéis definidos, capacidade de revisão humana), documentação auditável e comunicação compreensível ao titular de dados, articulando obrigações legais de proteção de dados com expectativas sociais de justiça e não discriminação \citep{ICO2022}.

No campo técnico, uma frente que tem recebido atenção crescente é a interpretabilidade mecanicista, cujo foco é identificar mecanismos internos responsáveis pelo comportamento do modelo, conectando componentes computacionais (como neurônios, circuitos e submódulos) aos resultados observados. Em vez de apenas descrever correlações de entrada e saída, essa abordagem busca explicações estruturais sobre \textit{como} e \textit{por que} a decisão foi produzida, aproximando a análise de noções de explicação causal e dedutiva no nível do mecanismo \citep{Kastner2024}. Em paralelo, os métodos geométricos oferecem uma base matemática complementar para a interpretabilidade ao modelar representações em estruturas não euclidianas, como grafos e variedades, e ao explorar propriedades de invariância, equivariância e similaridade entre estados internos da rede. A literatura de \textit{geometric deep learning} mostra que essa perspectiva permite construir descrições mais estáveis da organização das representações, favorecendo análises comparativas entre camadas e entre classes de entrada \citep{Cao2020}.

Portanto, este trabalho tem como objetivo propor um método geométrico de interpretabilidade baseado em \textit{kernels}. A fundamentação teórica do método proposto apoia-se nos paradigmas clássicos da aprendizagem de máquina, partindo dos \textit{Reproducing Kernel Hilbert Spaces} (RKHS) \citep{Aronszajn1950}, passando pelo aprendizado baseado em \textit{kernels} \citep{Scholkopf2001}, até alcançar a aprendizagem geométrica \citep{Cao2020} e a interpretabilidade mecanicista \citep{Kastner2024}. O método associa a cada entrada um conjunto de fontes de informação extraídas do próprio modelo durante o processamento dessa entrada, compara cada fonte entre as entradas por meio de um \textit{kernel} positivamente definido e compõe as matrizes de Gram resultantes pelo produto de Hadamard, que realiza o produto tensorial dos RKHS correspondentes. Dessa composição resulta um espaço de representação cujas distâncias e cuja distribuição de energia espectral são interpretadas como afirmações sobre o comportamento do modelo, isto é, sobre qual decisão ele toma, em quais casos ele erra e quais categorias ele confunde. Veja que o método é definido pela forma dessa construção, e não pelas fontes ou pelos \textit{kernels} escolhidos, de modo que ele é um método \textit{post-hoc} e modelo-agnóstico, aplicável a diferentes modelos e fontes de informação. Além disso, a distância induzida permite explicar a decisão sobre uma amostra por contraste com um contrafactual, a amostra mais próxima que o modelo decide de outra forma, indicando quais das diferenças entre as duas importam para o modelo.

O método é denominado \textbf{Fannar}, palavra islandesa ligada a \textit{fönn}, o monte de neve acumulado pelo vento (\textit{snow drift}). O nome descreve a ideia central do método. O vento não pode ser observado diretamente, mas a forma dos montes de neve que ele deposita registra a sua direção e a sua intensidade, e as camadas sucessivas de neve registram as tempestades na ordem em que ocorreram. Da mesma forma, o processo interno de decisão de um modelo não é observável diretamente, mas as amostras que ele processa se acumulam em um espaço de representação cuja forma registra como o modelo as trata, e a sequência de blocos registra em que profundidade cada distinção se forma. O método não explica o modelo por uma regra que lhe é imposta, e sim lendo a geometria que o próprio modelo deixa nos dados.

As contribuições deste trabalho são de duas naturezas, e é importante distingui-las. A primeira, e principal, é teórica. O tensor de similaridade em RKHS e a sua realização amostral pelo produto de Hadamard (Teorema \ref{thm:tensor_similarity} e Corolário \ref{cor:hadamard_form}), o pipeline núcleo-transformação com a sua classe de transformações admissíveis (Definições \ref{def:kernel_pipeline} e \ref{def:T_admissible}), a pseudométrica e a geometria relativa induzidas pela matriz de Gram (Seção \ref{sec:representation_geometry}) e a decomposição da energia espectral em uma parte retida e uma parte ortogonal (Subseção \ref{subsec:spectral}) constituem, em conjunto, um arcabouço formal para construir espaços de representação a partir de fontes heterogêneas de um mesmo modelo. Veja que esses resultados são demonstrados, e não medidos: eles valem para quaisquer fontes e quaisquer \textit{kernels} positivamente definidos, e não dependem dos experimentos para serem verdadeiros. A segunda contribuição é empírica. Ela compreende uma instanciação do método para redes neurais, com fontes extraídas do estado, do gradiente e dos parâmetros de cada bloco, uma segunda instanciação para árvores de decisão, que demonstra o caráter modelo-agnóstico da construção, e um protocolo de avaliação que repete cada leitura sobre onze bases de dados, três arquiteturas e várias sementes, sempre contra uma referência de nulidade construída para ela. Os experimentos, portanto, avaliam as escolhas dessa instanciação, e não a validade da construção, que é garantida pela teoria. Por fim, o método e o código dos experimentos são disponibilizados como a biblioteca de código aberto \texttt{fannar}\footnote{\url{https://github.com/FernandoCsm-Knight/Fannar}}.

